\section{Money Sums}
\label{sec:cses-money-sums}

\problemstatement{You have $n$ coins with given values. Determine all the
distinct sums you can form using \emph{some subset} of the coins (each
coin used at most once). Output how many distinct sums are possible and
list them.}

\begin{patternbox}
\emph{``Which totals are achievable from a subset''} is the boolean
\textbf{subset-sum} DP. Let $\textit{dp}[s]$ be true if some subset sums
to exactly $s$; each coin either joins a subset or not, a $0/1$ choice.
\end{patternbox}

\subsection*{State and transition}

$\textit{dp}[s]$ = whether sum $s$ is reachable using a subset of the
coins seen so far, with $\textit{dp}[0]=$ true (the empty subset). Adding
a coin of value $c$ makes any previously reachable $s$ also reach $s+c$.
Using a 1D array we iterate $s$ \textbf{downward} (0/1 knapsack rule) so a
coin is not reused:
\[
  \textit{dp}[s]\gets \textit{dp}[s]\ \lor\ \textit{dp}[s-c].
\]

\begin{ideabox}[Reachability, Not Counting]
Money Sums only asks \emph{whether} a sum is attainable, so the DP stores
booleans and combines with OR. The total possible sum is
$\sum c_i\le10^5$, so a single boolean array over $[0,\text{total}]$
suffices, and the answer is every index set to true.
\end{ideabox}

\subsection*{Algorithm}

\begin{enumerate}
  \item Let $S=\sum c_i$. Set $\textit{dp}[0]=$ true.
  \item For each coin $c$: for $s=S$ down to $c$,
        $\textit{dp}[s]\mathrel{|}=\textit{dp}[s-c]$.
  \item Collect all $s\ge1$ with $\textit{dp}[s]$ true; their count and
        list are the answer.
\end{enumerate}

\subsection*{Dry run}

\dryrun{Coins $\{4,2,2\}$; total $S=8$.}

\begin{itemize}
  \item Start $\textit{dp}[0]=$T. After coin $4$: $\{0,4\}$.
  \item After a $2$: $\{0,2,4,6\}$. After the second $2$:
        $\{0,2,4,6,8\}$.
  \item Non-zero reachable sums: $2,4,6,8$ -- four distinct sums.
\end{itemize}

\subsection*{C++ implementation}

\begin{lstlisting}
#include <bits/stdc++.h>
using namespace std;

int main() {
    int n;
    cin >> n;
    vector<int> coins(n);
    int total = 0;
    for (int& c : coins) { cin >> c; total += c; }

    vector<char> dp(total + 1, 0);
    dp[0] = 1;
    for (int c : coins)
        for (int s = total; s >= c; s--)           // downward: 0/1
            if (dp[s - c]) dp[s] = 1;

    vector<int> sums;
    for (int s = 1; s <= total; s++) if (dp[s]) sums.push_back(s);

    cout << sums.size() << "\n";
    for (int s : sums) cout << s << " ";
    cout << "\n";
    return 0;
}
\end{lstlisting}

\begin{complexitybox}
\textbf{Time Complexity.} $O(n\cdot S)$ where $S=\sum c_i$.
\textbf{Space Complexity.} $O(S)$ for the boolean array.
\end{complexitybox}

\begin{takeawaybox}
Subset-sum reachability is 0/1 knapsack with booleans and OR. The downward
loop again enforces ``each coin once.'' This reachability array is the
engine behind Two Sets II (next), where we \emph{count} the subsets rather
than merely test existence.
\end{takeawaybox}
